% ---
% id: transmitter/thm:certificate-sufficiency
% kind: thm
% title: "Checkability and sufficiency of the source-replacement certificate"
% summary: >-
%   In any form (pure projector, sector blocks, ensemble presentation) the
%   certificate's algebraic conditions are finitely many exact decisions
%   making Alice's measurement reproduce the virtual ensemble. Under retention
%   and sequential rounds she may measure before or after any attack: source
%   replacement.
% prominence: main
% status: implemented
% uses: [def:srcert, def:source-spec, lem:ensemble-presentation]
% mentions: [rem:xreg, fact:sr, rem:certificate-forms]
% assumes: [A-SEQ, A-SRC]
% refs: []
% provenance: newly-derived
% implements: certiqs_key.transmitter.certificate:verify_certificate
% lean:
% notes: "fact:sr is a mention and not a use: the proof of (iii)(b) names its source state only to identify the prepare-and-measure source state, the equality itself coming from (SR2). The implication (SR1) => tau_A in S_=(A) of (i) is used by no claim; req:marginal-decided is the rule that the marginal is decided directly. Also realised by certiqs_key.transmitter.certificate:consistency_criterion."
% ---

The theorem licenses source replacement at the boundary of this module.
Recall the conditions of Definition~\ref{def:srcert}: (SR1)
$\tr_{\Avirt}\Psi=\tau_A$, (SR2)
$\tr_A[(M_x\otimes\idon{\Avirt})\Psi]=p_x\rvirt_x$, and the attestation (SR3)
that nothing but Alice's measurement acts on $A$; an \tref{ensemble presentation}
has instead the decided conditions (E1) to (E3), positive declared blocks,
normalised weights and per-sector sums that round nowhere. Part (i) is about
declared data before its conditions are known to hold, and says that the
conditions other than the attestation (SR3) can be decided. Part (ii) relates
a certificate to Alice's measurement and is stated for any Hermitian family,
since its converse concerns a family that fails the POVM condition on the
support of $\tau_A$, that is, $\Pi(\sum_xM_x)\Pi\ne\Pi$ for the projector
$\Pi$ onto that support, which no \tref{source specification} declares, a
mismatch the \tref{pure form} detects (Remark~\ref{rem:certificate-forms}). Part
(iii) shows that her measurement may be moved across any attack that leaves
her retained system and her outcome register alone and starts from an
adversary register in a product state with the source. Its identities hold
for every channel of that form; the attestation (SR3) of
Assumption~\ref{ass:src} and Assumption~\ref{ass:seq} are what make the
attack of a round such a channel, and Remark~\ref{rem:xreg} shows that the
restriction on the attack is necessary and not a weakening of the intended
reading.

\begin{theorem}
\label{thm:certificate-sufficiency}
Let a \tref{source specification} (Definition~\ref{def:source-spec}) carry a
\tref{source-replacement certificate} $\Psi$ for $(A,\mathbb{T},\tau_A,\{M_x\})$
against the \tref{virtual ensemble} $\{p_x,\rvirt_x\}$, in any form of
Definition~\ref{def:srcert}. In (ii) and (iii) and their proofs an ensemble
presentation is read as the block-form certificate it defines
(Lemma~\ref{lem:ensemble-presentation}\,(i)), so that ``either form'' below
covers it.
\begin{enumerate}[label=(\roman*),leftmargin=2.2em]
\item \emph{(checkability)} the conditions on $\Psi$ of its declared form
--- in pure form Hermiticity, $\Psi^2=\Psi$ and $\tr\Psi=1$; in \tref{block form}
Hermiticity and positive semidefiniteness of every block and
$\sum_t\tr\Psi_t=1$ --- together with (SR1) and (SR2), are finitely many
exact identities and exact decisions on the declared data, decided by
finitely many arithmetic operations. In block form (SR1) and (SR2) are decided
as the block identities \eqref{eq:srblock} together with the block diagonality
of every weighted virtual member $\sigma^{\mathrm v}_x$ along $(S_t)_t$,
decided entry by entry, which by Definition~\ref{def:srcert} is their
conjunction; \eqref{eq:srblock} alone decides (SR2) only against the pinched
members $\bigoplus_t\Pi^{\mathrm v}_t\sigma^{\mathrm v}_x\Pi^{\mathrm v}_t$.
For a certificate given by an ensemble
presentation the decided conditions are (E1) to (E3) of
Definition~\ref{def:srcert}, equally finite and exact, and the conditions
of the block form hold of the defined certificate by
Lemma~\ref{lem:ensemble-presentation}\,(i) without being decided. (SR3) is
an attestation and is decided by no function of the declared data. (SR1)
moreover implies $\tau_A\in\Seq{A}$, $\tr_{\Avirt}\Psi$ being positive
semidefinite with trace $\tr\Psi=1$. In an ensemble presentation $\tau_A$ is
not declared: it is formed as the transposed sum of (E3), and that it is a
state is Lemma~\ref{lem:ensemble-presentation}\,(i).
\item \emph{(consistency)} more generally than in a source specification,
let $\{M_x\}$ be any family of Hermitian operators on $A$, presented as block
lists against the \tref{sector structure}, of which $\sum_xM_x=\idon{A}$ is not
assumed; let $\Psi$ meet the conditions on $\Psi$ of its form listed in (i)
and (SR1); and let $p_x\rvirt_x$ stand for
$\tr_A[(M_x\otimes\idon{\Avirt})\Psi]$, so that (SR2) holds. Writing
$U=\operatorname{supp}\tau_A$ with projector $\Pi_U$, the certificate's
signal-side marginal and the virtual average signal state agree,
\begin{equation}\label{eq:srconsist}
  \sum_xp_x\rvirt_x=\tr_A\Psi ,
\end{equation}
\emph{if} $\Pi_U\big(\sum_xM_x\big)\Pi_U=\Pi_U$, in either form; in
particular they agree for the certificate of a source specification, whose
$\{M_x\}$ is a POVM by Definition~\ref{def:source-spec}. The converse,
\emph{only if}, holds in pure form, and in block form whenever every non-zero
block $\Psi_t$ is a positive multiple of a rank-one projector; it fails for a
general positive block.
\item \emph{(sufficiency at the module boundary)} Assume (SR1), (SR2), the
attestation (SR3) of Assumption~\ref{ass:src}, and
Assumption~\ref{ass:seq}; in block form (SR2) is the identity of operators on
$\Avirt$ against the weighted virtual members $p_x\rvirt_x$ themselves, and
not against their diagonal blocks alone: for a declared block form (i) decides
it with their block diagonality, and for an ensemble presentation it is
proved by Lemma~\ref{lem:ensemble-presentation}\,(i). Let $\pcNmsm$ be Alice's measurement map, from
her retained system $A$ to the classical outcome register $X$,
\begin{equation}\label{eq:Nmsm}
  \pcNmsm(\xi)\coloneqq\sum_{x\in\Xset}\tr[M_x\,\xi]\;\dyad{x}_{X},
  \qquad \xi\in\Lin{A} ,
\end{equation}
and let $\mathcal{E}$ be \emph{any} channel from $\Avirt E$ to $E'$, where
$E$ is the adversary's initial register and $E'$ collects her output
registers including whatever is delivered to the receiver. Thus
$\mathcal{E}$ models the adversary's attack and the remainder of the round
\emph{on the \tref[virtual transmitted system]{transmitted system} and the adversary's registers}, and it
acts on neither $A$ nor $X$: it is not permitted to depend on Alice's
outcome. Write
$\rho_{X\Avirt}\coloneqq\sum_{x\in\Xset}p_x\,\dyad{x}_X\otimes\rvirt_x$ for
the virtual prepare-and-measure source state. Then:
\begin{enumerate}[label=(\alph*),leftmargin=2.2em]
\item \emph{(the interchange)}
\begin{equation}\label{eq:srmain}
  (\idmap_{X}\otimes\mathcal{E})\circ(\pcNmsm\otimes\idmap_{\Avirt E})
  \;=\;\pcNmsm\otimes\mathcal{E}\;=\;
  (\pcNmsm\otimes\idmap_{E'})\circ(\idmap_{A}\otimes\mathcal{E})
\end{equation}
as linear maps on $\Lin{A\Avirt E}$, hence in particular on
$\Psi\otimes\sigma$ for every state $\sigma$ on $E$: Alice's measurement may
be performed before the attack or after it, indifferently. This uses no
property of the certificate, only that the two maps act on complementary
registers;
\item \emph{(the source identification)} for every state $\sigma$ on $E$,
\begin{equation}\label{eq:srident}
  (\pcNmsm\otimes\mathcal{E})\big(\Psi\otimes\sigma\big)
  \;=\;(\idmap_{X}\otimes\mathcal{E})\big(\rho_{X\Avirt}\otimes\sigma\big) ,
\end{equation}
and both sides are states on $XE'$. Together with (a): on every register
other than $A$, the declared entanglement-based description and the
virtual prepare-and-measure description produce the same state, against
every attack of this form. The states $\rvirt_x$ in $\rho_{X\Avirt}$ are the
declared virtual states, in block form as in every form, and not their
pinching along $(S_t)_t$.
\end{enumerate}
\end{enumerate}
\end{theorem}

\begin{proof}
\emph{Idea.} Part (i) counts the operations each condition needs and
observes that they are formed exactly. Part (ii) reduces the consistency
identity to the vanishing of the compression $\Pi_UY\Pi_U$ of the defect
$Y=\sum_xM_x-\idon{A}$, using that $\Psi$ lives on the support of $\tau_A$.
Part (iii) is the interchange law for a tensor product of maps on
complementary registers, followed by (SR2).

(i) In pure form the conditions on $\Psi$ compare $(d_A\dim\Hs{\Avirt})^2$
entries each --- $\Psi=\Psi^{\dagger}$ and $\Psi^2=\Psi$ --- and one
accumulated trace; in block form they are, per block, one comparison of
$(d_{A_t}\dim\Hs{S_t})^2$ entries for Hermiticity and one exact decision of
positive semidefiniteness --- a symmetric elimination over the Gaussian
rationals in which a negative pivot refuses the block and a zero pivot
requires its whole row and column to vanish, the block being refused
otherwise, which decides that every principal minor is non-negative --- and
one accumulated trace over all blocks; in an ensemble presentation the same
decision per declared block, the traces of (E2) and the entrywise sums of
(E3); in pure form (SR1) compares $d_A^2$ entries after a partial trace and
(SR2) compares $|\Xset|$ matrices of size $\dim\Hs{\Avirt}$, each obtained
from $\Psi$ by a multiplication and a partial trace; in block form
\eqref{eq:srblock} compares, per sector, $d_{A_t}^2$ entries and $|\Xset|$
matrices of size $\dim\Hs{S_t}$, each obtained from $\Psi_t$ in the same way,
and the block diagonality compares with zero the entries of every
$\sigma^{\mathrm v}_x$ outside its diagonal blocks, at most
$|\Xset|(\dim\Hs{\Avirt})^2$ of them. By Definition~\ref{def:srcert} the block
identities and the block diagonality together are (SR1) and (SR2) of $\Psi$:
the partial traces of $\Psi$ are the direct sums of those of its blocks, and
the one of (SR2) is block diagonal along $(S_t)_t$, so it equals
$\sigma^{\mathrm v}_x$ exactly when $\sigma^{\mathrm v}_x$ has no entry
outside its diagonal blocks and those blocks satisfy the second identity of
\eqref{eq:srblock}. Each is a finite comparison, and each
is \emph{formed as well as taken} in exact arithmetic on the declared
dyadic entries (the number model of Definition~\ref{def:source-spec}): the
products entering $\Psi^2$ and (SR2), and the additions entering the
partial traces and the accumulated trace, are exact, so what is compared is
the quantity written and not a rounded image of it. That the operands are
exactly represented does not on its own deliver that; the exact formation
does. For the implication: $\tr_{\Avirt}\Psi$ is positive semidefinite as
the partial trace of a positive semidefinite operator --- in pure form
$\Psi$ is a projector, hence $\Psi=\Psi^{\dagger}\Psi\succeq0$; in block
form $\Psi$ is a direct sum of positive semidefinite blocks --- and its
trace is $\tr\Psi=1$; (SR1) is a matrix identity, so $\tau_A\in\Seq{A}$.
(SR3) is a statement about the apparatus, not about matrices, and no matrix
function decides it; it is the attestation of Assumption~\ref{ass:src}.

(ii) Put $Y\coloneqq\sum_xM_x-\idon{A}$, which is Hermitian, each $M_x$
being Hermitian by hypothesis and a sum of Hermitian matrices less
$\idon{A}$ being Hermitian; $Y$ is block diagonal along the sector structure,
every $M_x$ being presented as a block list. Summing (SR2) over $x$ gives
\begin{equation*}
  \tr_A\!\Big[\Big(\big(\textstyle\sum_xM_x\big)\otimes\idon{}\Big)\Psi\Big]
  =\sum_xp_x\rvirt_x ,
\end{equation*}
and the left-hand side is $\tr_A[(Y\otimes\idon{})\Psi]+\tr_A\Psi$ by
linearity of the partial trace, so \eqref{eq:srconsist} holds if and only
if $\tr_A[(Y\otimes\idon{})\Psi]=0$.

\emph{If, in either form.} $\Psi\succeq0$ in either form. Let $\ket{\phi}$
be a unit vector of $\Hs{A}$ orthogonal to $U$ and put
$P\coloneqq\dyad{\phi}\otimes\idon{\Avirt}$. By (SR1),
$\tr[P\Psi]=\bra{\phi}\tau_A\ket{\phi}=0$, so
$\Psi^{1/2}P\Psi^{1/2}$ is a positive semidefinite operator of trace zero,
hence zero, hence $P\Psi^{1/2}=0$ and $P\Psi=0$. Summing over an orthonormal
basis of the orthogonal complement of $U$ gives
$\Psi=(\Pi_U\otimes\idon{})\Psi(\Pi_U\otimes\idon{})$, and then
\begin{equation*}
  \tr_A\!\big[(Y\otimes\idon{})\,\Psi\big]
  =\tr_A\!\big[(Y\Pi_U\otimes\idon{})\,\Psi\,(\Pi_U\otimes\idon{})\big]
  =\tr_A\!\big[(\Pi_UY\Pi_U\otimes\idon{})\,\Psi\big] ,
\end{equation*}
the partial trace over $A$ being cyclic under operators that act on $A$
alone. This vanishes when $\Pi_UY\Pi_U=0$, i.e.\ when
$\Pi_U(\sum_xM_x)\Pi_U=\Pi_U$. If $\{M_x\}$ is a POVM then $Y=0$ and the
condition holds outright, and Definition~\ref{def:source-spec} delivers
exactly that: $\sum_xM^{(t)}_x=\idon{A_t}$ for every $t$ as an identity of
the declared numbers, its declared-order sum decided exactly.

\emph{Only if, in pure form.} $\Psi=\dyad{\psi}$ for a unit vector
$\ket{\psi}$; take a Schmidt decomposition
$\ket{\psi}=\sum_i\sqrt{\mu_i}\ket{e_i}_A\ket{f_i}_{\Avirt}$ with every
$\mu_i>0$ and $\{\ket{e_i}\}$, $\{\ket{f_i}\}$ orthonormal. By (SR1),
$\tau_A=\tr_{\Avirt}\Psi=\sum_i\mu_i\dyad{e_i}$, which is a spectral
decomposition, so $\{\ket{e_i}\}$ is an orthonormal basis of
$U=\operatorname{supp}\tau_A$. Then
\begin{equation*}
  \tr_A\!\big[(Y\otimes\idon{})\,\Psi\big]
  =\sum_{i,j}\sqrt{\mu_i\mu_j}\,\bra{e_j}Y\ket{e_i}\,\ketbra{f_i}{f_j},
\end{equation*}
and since the $\ket{f_i}$ are orthonormal and every $\mu_i>0$, this
operator vanishes if and only if $\bra{e_j}Y\ket{e_i}=0$ for all $i,j$,
i.e.\ if and only if $\Pi_UY\Pi_U=0$.

\emph{Only if, in block form with rank-one blocks.} Let
$\Psi_t=w_t\dyad{\psi_t}$ with $w_t>0$ and $\ket{\psi_t}$ a unit vector of
$\Hs{A_tS_t}$ for every non-zero block. $Y=\bigoplus_tY_t$ is block diagonal
because every $M_x$ is, so
$\tr_A[(Y\otimes\idon{})\Psi]=\bigoplus_tw_t\,
\tr_{A_t}[(Y_t\otimes\idon{S_t})\dyad{\psi_t}]$, a direct sum along
$(S_t)_t$, which vanishes if and only if every summand does. The pure-form
argument, applied to $\ket{\psi_t}$ with $\tau^{(t)}_A/w_t$ in the place of
$\tau_A$, shows that the $t$-th summand vanishes if and only if
$\Pi_{U,t}Y_t\Pi_{U,t}=0$, where $\Pi_{U,t}$ projects onto
$\operatorname{supp}\tau^{(t)}_A$; and
$\Pi_U=\bigoplus_t\Pi_{U,t}$, so the conjunction over $t$ is
$\Pi_UY\Pi_U=0$.

\emph{The converse fails for a general positive block.} Take one sector,
$\Psi=\tau_A\otimes\varsigma$ for a state $\varsigma$ on $\Avirt$, and any
Hermitian $Y$ with $\tr[Y\tau_A]=0$ and $\Pi_UY\Pi_U\ne0$; such a $Y$
exists whenever $\dim U\ge2$, and scaled so that $\idon{A}+Y\succeq0$ it is
realised by the one-letter family $M=\idon{A}+Y$, Hermitian and positive
semidefinite but not a POVM, which the hypothesis of (ii) admits, with
$p\rvirt=\tr_A[(M\otimes\idon{})\Psi]$ as in the statement. The block
$\Psi$ is positive semidefinite of unit trace and satisfies (SR1). Then
$\tr_A[(Y\otimes\idon{})\Psi]=\tr[Y\tau_A]\,\varsigma=0$.

(iii)(a) \emph{The interchange.} The registers do the work, and no
numerical hypothesis enters. The input operator lives on $A\Avirt E$. By
hypothesis $\pcNmsm$ is a linear map out of $A$ alone and $\mathcal{E}$ a
channel out of $\Avirt E$ alone, and $A$ and $\Avirt E$ are complementary
tensor factors of the input; the outputs $X$ and $E'$ are likewise
distinct registers. Hence both composites equal the tensor product
$\pcNmsm\otimes\mathcal{E}$ of the two maps, which is \eqref{eq:srmain};
this is the elementary interchange law for a tensor product of linear maps
and needs nothing about the particular $\pcNmsm$, $\mathcal{E}$ or input
operator. That $\mathcal{E}$ does not act on $A$ is (SR3); that it does not
act on $X$ is a hypothesis of the statement, and Remark~\ref{rem:xreg}
shows it cannot be dropped.

(iii)(b) \emph{The source identification.} Applying
$\pcNmsm\otimes\idmap_{\Avirt E}$ to $\Psi\otimes\sigma$ gives, by
\eqref{eq:Nmsm} and (SR2),
\begin{equation*}
  \sum_x\dyad{x}_X\otimes
  \tr_A\!\big[(M_x\otimes\idon{\Avirt})\,\Psi\big]\otimes\sigma
  =\sum_xp_x\,\dyad{x}_X\otimes\rvirt_x\otimes\sigma
  =\rho_{X\Avirt}\otimes\sigma ,
\end{equation*}
the right-hand side being exactly the classical-quantum source state from
which Fact~\ref{fact:sr} starts, under the symbol map stated with that
fact. Applying the linear map $\idmap_X\otimes\mathcal{E}$ to both
sides gives \eqref{eq:srident}, the left-hand side being
$(\pcNmsm\otimes\mathcal{E})(\Psi\otimes\sigma)$ by (a). That both sides
are states: $\{M_x\}$ is a POVM by Definition~\ref{def:source-spec}, so
$\pcNmsm$ is a channel; $\Psi$ is a state, being positive semidefinite of
unit trace in either form; and $\mathcal{E}$ is a channel by hypothesis.
Nothing in (iii) reads $\Psi^2=\Psi$.
\end{proof}
